\chapter{HiED Architecture and Stage-Wise Diagnostic Workflow}
\label{ch:architecture}

This chapter explains how HiED transforms one fixed psychiatric interview transcript into separately recorded diagnostic outputs. It first defines the task input, diagnostic scope, and possible outputs. It then presents the two-path architecture, the primary-diagnosis selection policies, the Single LLM comparison contract, and a complete worked example.

\section{Task Input, Diagnostic Scope, and Recorded Outputs}
\label{sec:input-diagnostic-scope}

\subsection{Task Input}

HiED processes one speaker-labeled Chinese psychiatric interview transcript at a time. The transcript contains an ordered sequence of clinician and patient turns. All compared methods receive the same fixed transcript and cannot ask the patient follow-up questions.

When similar-case retrieval is enabled, labeled examples from the training source are provided only to the diagnosis path. These examples may guide the model, but they are not evidence about the current patient.

The criterion-checking path also receives the configured diagnostic criteria used in this study. These criteria are fixed study rules rather than information collected from the current patient.

\subsection{Configured Diagnostic Scope}

HiED uses fourteen configured ICD-10 diagnostic categories. These categories define the complete diagnostic space that the system may rank, check, or select. A diagnosis outside this space cannot be produced by HiED.

The Criterion Checkers evaluate all fourteen categories for every transcript. F41 and F43 subtypes remain separate during criterion checking, but they are mapped to broader parent labels for evaluation. The complete operational compatibility rules are reported in Appendix~\ref{app:configured-profile}.

\begin{table}[htbp]
\centering
\caption{Configured ICD-10 diagnostic categories and their scoring parent labels.}
\label{tab:configured-profile-summary}
\small
\begin{tabularx}{\textwidth}{|p{3.0cm}|p{2.2cm}|X|}
\hline
ICD-10 code & Scoring parent & Diagnostic category\\
\hline
F20 & F20 & Schizophrenia\\
\hline
F31 & F31 & Bipolar affective disorder\\
\hline
F32 & F32 & Depressive episode\\
\hline
F39 & F39 & Unspecified mood disorder\\
\hline
F41.0 & F41 & Panic disorder\\
\hline
F41.1 & F41 & Generalized anxiety disorder\\
\hline
F41.2 & F41 & Mixed anxiety and depressive disorder\\
\hline
F42 & F42 & Obsessive-compulsive disorder\\
\hline
F43.1 & F43 & Post-traumatic stress disorder\\
\hline
F43.2 & F43 & Adjustment disorders\\
\hline
F45 & F45 & Somatoform disorders\\
\hline
F51 & F51 & Nonorganic sleep disorders\\
\hline
F98 & F98 & Other behavioral and emotional disorders\\
\hline
Z71 & Z71 & Counseling without a specific disorder\\
\hline
\end{tabularx}

\vspace{0.4em}
\begin{minipage}{0.96\textwidth}
\footnotesize
\textit{Note.} F41.0, F41.1, and F41.2 are mapped to F41 for evaluation. F43.1 and F43.2 are mapped to F43. The evaluation-only \emph{Others} label is not part of the configured output space and cannot be emitted by HiED.
\end{minipage}
\end{table}

\subsection{Recorded Output Types}
\label{sec:diagnostic-output-roles}

HiED produces more than one final diagnosis. For each transcript, it keeps several outputs with different roles. Table~\ref{tab:hied-output-space} summarizes these outputs. A candidate diagnosis, a compatible diagnosis, a comorbid diagnosis, and a committed primary diagnosis should not be treated as the same type of output.

\begin{table}[htbp]
\centering
\caption{Recorded outputs produced by HiED for one transcript.}
\label{tab:hied-output-space}
\label{tab:diagnostic-output-roles}
\footnotesize
\begin{tabularx}{\textwidth}{|>{\raggedright\arraybackslash}p{3.4cm}|>{\centering\arraybackslash}p{2.3cm}|>{\raggedright\arraybackslash}X|}
\hline
\textbf{Recorded output} & \textbf{Possible amount} & \textbf{Role}\\
\hline
Ranked candidate list & Up to five & Ordered differential diagnoses proposed by the diagnosis path.\\
\hline
Proposed primary diagnosis & One & The Diagnostician's direct first choice before finalization.\\
\hline
Optional comorbid diagnosis & Zero or one & An additional diagnosis proposed as being present with the primary diagnosis.\\
\hline
Criterion states & All configured criteria & Whether each criterion is \texttt{met}, \texttt{not\_met}, or \texttt{insufficient\_\allowbreak{}evidence}.\\
\hline
Evidence notes & When relevant & Transcript information linked by the checker to a criterion.\\
\hline
Missing-information record & When evidence is insufficient & Information that is not available in the fixed transcript.\\
\hline
Criterion-compatible set & A subset of the fourteen categories & Diagnoses that pass their study-specific compatibility rules.\\
\hline
Committed primary diagnosis & One & The final primary diagnosis recorded under the selected finalization policy.\\
\hline
\end{tabularx}
\end{table}

\section{HiED Two-Path Architecture}
\label{sec:hied-two-path-architecture}

HiED uses two connected paths to examine the same fixed transcript. The diagnosis path forms and ranks possible diagnoses. The criterion-checking path records criterion-level evidence for every configured diagnosis. Their outputs remain separate until the primary-selection step.

The two paths are separate in their data flow and output roles. This does not mean that their errors are statistically independent. Both use the same transcript, and their model-based components may show similar behavior.

Figure~\ref{fig:hied-two-path-architecture} summarizes the workflow. Similar cases, when used, enter only the diagnosis path. The criterion-checking path receives the fixed transcript and the configured criteria for all fourteen diagnoses. Direct-Answer and Nominate-then-Select use the recorded outputs differently when they commit the final primary diagnosis.

\begin{figure}[htbp]
\centering
\includegraphics[width=1.0\textwidth]{fig_hied_two_path_architecture.pdf}
\caption{HiED two-path architecture and recorded diagnostic outputs. The fixed transcript is processed by a diagnosis path and a criterion-checking path. Similar cases, when used, are provided only to the diagnosis path. The diagnosis path records ranked candidates and a proposed primary diagnosis. The criterion-checking path records criterion states, evidence notes, missing information, and a criterion-compatible set. Direct-Answer keeps the proposed primary, whereas Nominate-then-Select uses the ranking and compatible set to reselect the committed primary. All intermediate records remain available for review.}
\label{fig:hied-two-path-architecture}
\label{fig:single-vs-hied-architecture}
\label{fig:pipeline-thesis-rewrite}
\end{figure}
\FloatBarrier

\section{Diagnosis Path: Candidate Formation and Primary Proposal}
\label{sec:retrieval}
\label{sec:diagnosis-path}

The diagnosis path forms a differential diagnosis from the fixed transcript. Its main outputs are a ranked candidate list, a proposed primary diagnosis, and an optional comorbid diagnosis.

\subsection{Optional Similar-Case Retrieval}

When retrieval is enabled, similar labeled cases are provided only to the diagnosis path. They may guide candidate formation and the proposed primary diagnosis, but they do not change the current transcript or the input to the Criterion Checkers. The retrieval policies and their evaluation are described in Chapter~\ref{ch:experimental}.

\subsection{Diagnostician Outputs}
\label{sec:diagnostician-outputs}

The Diagnostician receives the speaker-labeled transcript, the configured diagnostic scope, and any retrieved examples. It records:

\begin{enumerate}
    \item a ranked list of up to five candidate diagnoses;
    \item one proposed primary diagnosis selected from the leading candidates; and
    \item at most one optional comorbid diagnosis.
\end{enumerate}

The ranked list keeps possible alternatives visible even when they are not selected as primary. The proposed primary diagnosis is the Diagnostician's direct first choice. The optional comorbid diagnosis has a different role: it proposes that another disorder may be present at the same time. A diagnosis in the ranked list is therefore not automatically a comorbidity.

These three outputs come from the same Diagnostician and should not be treated as independent model opinions. They are kept separately because candidate availability, primary preference, and comorbidity are different diagnostic roles.

\section{Criterion-Checking Path: Criterion States and Compatibility}
\label{sec:criterion-checking-path}

The criterion-checking path examines the transcript at the diagnostic criterion level. It asks whether each configured diagnosis is compatible with the evidence in the fixed transcript. It does not rank the compatible diagnoses or decide which one should be primary.

\subsection{Diagnosis-Specific Criterion Checkers}
\label{sec:criterion-checkers}

Each of the fourteen configured diagnostic categories has its own Criterion Checker. Every checker receives the same speaker-labeled transcript together with the study-specific criteria for one diagnosis. The checkers do not receive the retrieved similar-case examples, and they do not use the Diagnostician's ranked candidate list.

All fourteen categories are checked for every case. A diagnosis may therefore enter the criterion-compatible set even when it was not ranked among the leading candidates.

For each criterion, the checker records one of three states. Table~\ref{tab:criterion-state-definitions} gives their meanings.

\begin{table}[htbp]
\centering
\caption{Criterion states recorded by the Criterion Checkers.}
\label{tab:criterion-state-definitions}
\small
\begin{tabularx}{\textwidth}{|>{\raggedright\arraybackslash}p{3.6cm}|>{\raggedright\arraybackslash}X|}
\hline
\textbf{Criterion state} & \textbf{Meaning in this study}\\
\hline
\texttt{met} & The transcript contains information that supports the criterion.\\
\hline
\texttt{not\_met} & The transcript contains information showing that the criterion is not met.\\
\hline
\texttt{insufficient\_\allowbreak{}evidence} & The transcript does not contain enough information to decide whether the criterion is met.\\
\hline
\end{tabularx}
\end{table}

An \texttt{insufficient\_\allowbreak{}evidence} state is not treated as negative evidence. It means that the current transcript does not provide enough information for either a \texttt{met} or a \texttt{not\_met} judgment. When relevant information is found, the checker also records a short evidence note for later review.

\subsection{Compatibility Auditor}
\label{sec:compatibility-auditor}

After the Criterion Checkers finish, a deterministic Compatibility Auditor applies the study-specific rule for each diagnosis. A rule may require particular core criteria, a minimum number of supported criteria, a required duration, or a combination of these conditions.

The auditor uses only the recorded criterion states. It does not reread the transcript or create new evidence. Diagnoses that pass their configured rules are placed in the criterion-compatible set. The complete rules are reported in Appendix~\ref{app:configured-profile}.

\subsection{Meaning of Criterion Compatibility}
\label{sec:criterion-compatibility-boundary}

A compatible diagnosis has passed the study-specific rule using the recorded criterion states. This does not show that it is clinically correct, mutually exclusive with other diagnoses, or the preferred primary diagnosis. Compatibility records whether one diagnosis fits its own rule; primary selection compares several possible diagnoses.

\section{Primary-Diagnosis Selection}
\label{sec:primary-selection}

The diagnosis and criterion-checking paths provide different views of the same transcript. The primary-selection step records which diagnosis is finally committed as the system output.

\subsection{Direct-Answer Policy}
\label{sec:direct-answer-policy}

Direct-Answer (DA) keeps the Diagnostician's proposed primary diagnosis as the committed primary diagnosis. The criterion records and criterion-compatible set remain available for review but do not replace that choice. DA may also retain one optional comorbid diagnosis.

\subsection{Nominate-then-Select Policy}
\label{sec:nominate-then-select-policy}

Nominate-then-Select (NtS) reselects the committed primary diagnosis using the ranked candidates and the criterion-compatible set. It returns one committed primary diagnosis and does not retain DA's optional comorbid output. The complete fallback order and tie-handling rules are reported in the experimental protocol and appendix.

\subsection{Shared Upstream Records and Comparison Boundary}
\label{sec:da-nts-relationship}
\label{sec:primary-selection-boundary}

DA and NtS reuse the same ranked candidates, proposed primary diagnosis, criterion states, and criterion-compatible set. Their difference is how these records are used to commit one primary diagnosis. Their committed Top-1 results can therefore be compared directly. Set-based measures also reflect their different output sizes because DA may retain a comorbid diagnosis whereas NtS returns one diagnosis.

The committed primary is the system's benchmark output. A mismatch with a gold label may reflect an unsuitable system choice, missing transcript information, or ambiguity in the dataset label. This thesis therefore reports benchmark agreement and disagreement rather than treating every mismatch as a confirmed clinical error.

\section{Single LLM Baseline and Comparison Contract}
\label{sec:compared-architectures}

\subsection{Single LLM Baseline}
\label{sec:single-architecture}

The Single LLM receives the same fixed transcript and optional similar-case examples. It returns a primary diagnosis and may emit additional diagnostic labels. These additional labels are not produced as the same ranked differential used by HiED, and the Single LLM does not produce a diagnosis-specific criterion record.

The Single LLM is therefore used as a direct diagnostic baseline. Its final diagnostic outputs can be compared with HiED, but its emitted labels cannot be treated as the same stage-wise records.

\subsection{Recorded Output Comparison}
\label{sec:recorded-output-differences}

Table~\ref{tab:single-hied-output-comparison} summarizes the outputs available from the two study architectures.

\begin{table}[htbp]
\centering
\caption{Recorded outputs available from the Single LLM and HiED study architectures.}
\label{tab:single-hied-output-comparison}
\footnotesize
\begin{tabularx}{\textwidth}{|>{\raggedright\arraybackslash}p{4.2cm}|>{\centering\arraybackslash}p{3.0cm}|>{\centering\arraybackslash}X|}
\hline
\textbf{Recorded output} & \textbf{Single LLM} & \textbf{HiED}\\
\hline
Committed primary diagnosis & Yes & Yes\\
\hline
Optional additional diagnosis & Possible under the Single format & Optional comorbid diagnosis\\
\hline
Genuine ranked differential diagnosis & No & Yes\\
\hline
Criterion state for each configured diagnosis & No & Yes\\
\hline
Evidence notes and missing-information record & No & Yes\\
\hline
Criterion-compatible diagnosis set & No & Yes\\
\hline
Separately observable diagnostic stages & Final emitted output only & Candidate generation, criterion checking, and primary selection\\
\hline
\end{tabularx}
\end{table}

The comparison concerns complete study architectures, not the isolated effect of agent count. The systems differ in workflow, model calls, intermediate outputs, and output contract.

\section{Complete Worked Example}
\label{sec:running-example}
\label{sec:worked-example}

\subsection{Example Purpose and Overview}

This section follows one complete constructed psychiatric interview through the Single LLM and HiED output contracts. The transcript is written for this thesis. It is not a patient record, a case copied from LingxiDiag, or an exported model trace. The example is used only to explain the data flow and output roles.

Figure~\ref{fig:worked-example-flow} provides an overview before the complete transcript and output tables. Only selected categories are shown in the figure, although the full HiED configuration checks all fourteen categories.

\begin{figure}[htbp]
\centering
\includegraphics[width=1.0\textwidth]{fig_worked_example_flow.pdf}
\caption{Illustrative HiED data flow for the complete constructed psychiatric interview. The diagnosis path ranks F41.1 above F32 and proposes F41.1 as primary. In the selected criterion records, F32 passes the configured compatibility rule, whereas the duration and associated symptoms required for F41.1 remain unclear. Direct-Answer keeps the proposed F41.1 primary diagnosis, while Nominate-then-Select commits F32 using the same upstream outputs. The example explains the output contract and does not represent a patient record, benchmark result, or clinically adjudicated diagnosis.}
\label{fig:worked-example-flow}
\end{figure}
\FloatBarrier

@@TRANSCRIPT_BLOCK@@

@@SINGLE_BLOCK@@

@@HIED_BLOCK@@

@@WHAT_BLOCK@@

\section{Chapter Summary}

HiED processes one fixed transcript within a fourteen-category diagnostic space. It separately records ranked candidates, criterion states, evidence notes, missing information, criterion compatibility, and the committed primary diagnosis.

The diagnosis path forms and ranks possible diagnoses and proposes a primary diagnosis. The criterion-checking path evaluates all configured categories and records which criteria are supported, contradicted, or unclear. A deterministic Compatibility Auditor then forms the criterion-compatible set.

DA and NtS use the same upstream records in different ways during primary selection. DA keeps the proposed primary diagnosis, whereas NtS uses the ranking and criterion-compatible set to make a new selection. The Single LLM provides a direct diagnostic baseline but does not provide the same stage-wise output contract.

The worked example showed how the same transcript can produce a ranked differential, criterion-level records, missing-information states, and different DA and NtS commitments. These records form the basis of the stage-wise evaluation defined in the following chapters.
